\section{Metodología experimental}
\label{sec:metodologia_experimental}

La metodología experimental combina optimización MILP, problemas auxiliares
de factibilidad, construcción de testigos y búsqueda exhaustiva restringida.
El propósito no es únicamente encontrar trayectorias con baja actividad,
sino determinar qué afirmaciones pueden considerarse mínimos globales y
cuáles deben presentarse únicamente como cotas.

Todos los candidatos utilizados en los resultados se reconstruyen mediante
la implementación funcional descrita en la
Sección~\ref{sec:arquitectura_implementacion}. De esta manera, la actividad
reportada no depende exclusivamente del valor objetivo entregado por CBC.


\subsection{Diseño experimental}
\label{subsec:diseno_experimental}

Se estudian los tamaños de palabra

\[
z \in \{4,8\}
\]

y los números de rondas

\[
R \in \{1,2,3\}.
\]

La combinación de ambos parámetros produce los seis experimentos de la
Tabla~\ref{tab:matriz_experimentos}.

\begin{table}[htbp]
    \centering
    \caption{Matriz de experimentos para Keccak reducido.}
    \label{tab:matriz_experimentos}
    \begin{tabular}{ccccc}
        \hline
        \textbf{Experimento}
        &
        \boldmath{$z$}
        &
        \textbf{Estado}
        &
        \textbf{S-boxes por ronda}
        &
        \textbf{Rondas}
        \\
        \hline
        E1 & 4 & 100 bits & 20 & 1 \\
        E2 & 4 & 100 bits & 20 & 2 \\
        E3 & 4 & 100 bits & 20 & 3 \\
        E4 & 8 & 200 bits & 40 & 1 \\
        E5 & 8 & 200 bits & 40 & 2 \\
        E6 & 8 & 200 bits & 40 & 3 \\
        \hline
    \end{tabular}
\end{table}

La correspondencia entre rangos de intentos y número de rondas se utiliza
únicamente para organizar estos escenarios, tal como se explicó en la
Subsección~\ref{subsec:intentos_rondas}.


\subsection{Variable de respuesta}
\label{subsec:variable_respuesta}

La variable principal es el número total de aplicaciones activas de
$\chi$ durante $R$ rondas:

\[
N_{\mathrm{act}}
=
\sum_{r=0}^{R-1}
\sum_{y=0}^{4}
\sum_{k=0}^{z-1}
a_{r,y,k}.
\]

La actividad de una ronda individual se representa mediante

\[
m_r
=
\sum_{y=0}^{4}
\sum_{k=0}^{z-1}
a_{r,y,k}.
\]

Por tanto,

\[
N_{\mathrm{act}}
=
m_0+m_1+\cdots+m_{R-1}.
\]

La notación

\[
2+2+9
\]

indica que una trayectoria activa dos S-boxes en la primera ronda, dos en
la segunda y nueve en la tercera.

El mínimo global para una configuración $(z,R)$ se denota por

\[
N_{\mathrm{act}}^{\min}(z,R).
\]

Cuando el valor exacto no puede certificarse, se reporta mediante un
intervalo

\[
L(z,R)
\leq
N_{\mathrm{act}}^{\min}(z,R)
\leq
U(z,R),
\]

donde $L(z,R)$ es una cota inferior demostrada y $U(z,R)$ es la actividad
de un testigo factible.


\subsection{Niveles de certificación}
\label{subsec:niveles_certificacion}

Los resultados se clasifican según el tipo de evidencia disponible. La
Tabla~\ref{tab:certificacion_resultados} resume las categorías utilizadas.

\begin{table}[htbp]
    \centering
    \caption{Niveles de certificación utilizados en los experimentos.}
    \label{tab:certificacion_resultados}
    \begin{tabular}{p{0.25\textwidth}p{0.67\textwidth}}
        \hline
        \textbf{Categoría} & \textbf{Interpretación} \\
        \hline

        Cota inferior
        &
        Valor por debajo del cual se ha demostrado que no puede existir
        una trayectoria válida.
        \\[3pt]

        Testigo factible
        &
        Pareja concreta de estados que satisface el modelo y alcanza una
        actividad determinada. Su valor proporciona una cota superior.
        \\[3pt]

        Óptimo global
        &
        La cota inferior coincide con la actividad de un testigo factible.
        No existe una solución mejor en el espacio completo del modelo.
        \\[3pt]

        Óptimo restringido
        &
        Mejor valor obtenido exhaustivamente dentro de una familia
        específica de trayectorias. No implica optimalidad global.
        \\
        \hline
    \end{tabular}
\end{table}

Una terminación por límite de tiempo no se interpreta como prueba de
optimalidad ni de infactibilidad.


\subsection{Procedimiento general}
\label{subsec:procedimiento_general}

Cada experimento sigue las siguientes etapas:

\begin{enumerate}
    \item seleccionar el tamaño de palabra $z$ y el número de rondas $R$;

    \item construir las dos ejecuciones del modelo MILP;

    \item imponer la condición de diferencia inicial no nula;

    \item incorporar las variables de actividad y la función objetivo;

    \item resolver el problema de minimización o un problema auxiliar de
    factibilidad;

    \item extraer los estados izquierdo y derecho cuando se obtiene una
    solución;

    \item reconstruir la trayectoria con la implementación funcional;

    \item calcular nuevamente las diferencias y los soportes activos;

    \item comparar la actividad reconstruida con la actividad obtenida por
    el modelo;

    \item clasificar el resultado según el nivel de certificación
    alcanzado.
\end{enumerate}

Los procedimientos concretos difieren según el número de rondas, debido al
crecimiento del espacio combinatorio.


\subsection{Certificación para una ronda}
\label{subsec:certificacion_una_ronda}

Para una ronda se impone

\[
\Delta A_0 \neq 0.
\]

Como las transformaciones $\theta$, $\rho$ y $\pi$ son biyectivas, una
diferencia inicial no nula no puede convertirse en una diferencia nula
antes de alcanzar la capa $\chi$:

\[
\Delta A_0 \neq 0
\Longrightarrow
\Delta B_0 \neq 0.
\]

En consecuencia, al menos una aplicación local de $\chi$ debe estar
activa:

\[
m_0 \geq 1.
\]

Esta propiedad proporciona una cota inferior teórica. Para completar la
certificación se busca una pareja concreta de estados cuya actividad sea

\[
m_0=1.
\]

Cuando el testigo es reconstruido correctamente, la coincidencia entre la
cota inferior y la cota superior establece el mínimo global.


\subsection{Certificación para dos rondas}
\label{subsec:certificacion_dos_rondas}

Para dos rondas se minimiza

\[
N_{\mathrm{act}}
=
m_0+m_1.
\]

La certificación combina un problema de factibilidad y la construcción de
un testigo.

En primer lugar, se añade la restricción

\[
m_0+m_1 \leq 3.
\]

Si CBC demuestra que este modelo es infactible, queda descartada cualquier
trayectoria con actividad total menor o igual que tres.

En segundo lugar, se busca y reconstruye una trayectoria con distribución

\[
m_0=2,
\qquad
m_1=2.
\]

Su actividad total es

\[
N_{\mathrm{act}}=4.
\]

La coincidencia entre la infactibilidad del problema con cota tres y la
existencia de un testigo de actividad cuatro permite certificar el mínimo
global para dos rondas.


\subsection{Estrategia para tres rondas}
\label{subsec:estrategia_tres_rondas}

Para tres rondas, la actividad total es

\[
N_{\mathrm{act}}
=
m_0+m_1+m_2.
\]

Toda trayectoria válida de tres rondas contiene en sus dos primeras rondas
una trayectoria válida de dos rondas. Además, como una ronda de Keccak es
biyectiva, la diferencia no puede desaparecer antes de la tercera
aplicación de $\chi$. Por tanto,

\[
m_2 \geq 1.
\]

Una vez certificado el mínimo de dos rondas, puede establecerse la cota

\[
N_{\mathrm{act}}^{\min}(z,3)
\geq
N_{\mathrm{act}}^{\min}(z,2)+1.
\]

Esta relación proporciona una cota inferior global para tres rondas.

Sin embargo, la resolución completa del modelo de tres rondas no permitió
cerrar la brecha de optimalidad con CBC. Por ello, la obtención de cotas
superiores se abordó mediante una búsqueda restringida a trayectorias con
estructura

\[
2+2+c,
\]

donde $c$ es la actividad inducida en la tercera ronda.


\subsection{Búsqueda exhaustiva restringida}
\label{subsec:busqueda_exhaustiva_restringida}

La búsqueda toma como punto de partida trayectorias óptimas de dos rondas
con distribución

\[
2+2.
\]

Una trayectoria se caracteriza mediante los soportes activos

\[
S_0,
S_1
\subseteq
\{0,\ldots,4\}
\times
\{0,\ldots,z-1\},
\]

con

\[
|S_0|=|S_1|=2.
\]

Después de aplicar los filtros de compatibilidad diferencial, se obtuvieron

\[
200
\]

trayectorias para $z=4$ y

\[
400
\]

trayectorias para $z=8$.

Estos conteos corresponden a las trayectorias admitidas por la rutina de
búsqueda implementada.


\subsubsection{Realizaciones concretas}

El soporte diferencial no determina de manera única los valores absolutos
de las entradas a $\chi$. Debido a la no linealidad de esta capa, diferentes
valores absolutos pueden producir distintas diferencias de salida.

Cada S-box local procesa cinco bits y admite

\[
2^5=32
\]

valores posibles. Como la segunda ronda contiene dos S-boxes activas, se
enumeran

\[
32^2=1024
\]

realizaciones por trayectoria.

Las posiciones cuya diferencia es nula se fijan a un representante común
en ambas ejecuciones. Esta elección evita enumerar valores equivalentes que
no modifican la diferencia local.

El tamaño total de las búsquedas se presenta en la
Tabla~\ref{tab:tamano_busqueda_restringida}.

\begin{table}[htbp]
    \centering
    \caption{Tamaño de la búsqueda exhaustiva restringida.}
    \label{tab:tamano_busqueda_restringida}
    \begin{tabular}{cccc}
        \hline
        \boldmath{$z$}
        &
        \textbf{Trayectorias}
        &
        \textbf{Realizaciones por trayectoria}
        &
        \textbf{Total}
        \\
        \hline
        4 & 200 & 1024 & 204\,800 \\
        8 & 400 & 1024 & 409\,600 \\
        \hline
    \end{tabular}
\end{table}


\subsubsection{Evaluación de la tercera ronda}

Para cada realización se propagan los estados hasta la entrada de $\chi$
en la tercera ronda. El soporte activo se calcula como

\[
S_2
=
\left\{
(y,k):
\bigvee_{x=0}^{4}
\Delta B_2[x,y,k]=1
\right\}.
\]

La actividad de la tercera ronda es

\[
c=|S_2|,
\]

y la actividad total del candidato es

\[
2+2+c.
\]

La búsqueda conserva los candidatos con menor valor de $c$ junto con la
información necesaria para reconstruir sus estados iniciales, sus
diferencias y sus soportes activos.


\subsection{Validación de los candidatos}
\label{subsec:validacion_candidatos}

Los mejores candidatos de la búsqueda restringida se validan mediante el
siguiente procedimiento:

\begin{enumerate}
    \item reconstruir los estados iniciales $A_0^L$ y $A_0^R$;

    \item ejecutar las tres rondas con la implementación de referencia;

    \item calcular las diferencias antes de cada capa $\chi$;

    \item reconstruir los soportes $S_0$, $S_1$ y $S_2$;

    \item comprobar que

    \[
    |S_0|=2,
    \qquad
    |S_1|=2;
    \]

    \item comprobar que $|S_2|$ coincida con el valor obtenido durante la
    enumeración;

    \item verificar que la actividad total sea

    \[
    2+2+|S_2|;
    \]

    \item fijar, cuando sea necesario, los estados y soportes dentro del
    modelo MILP para comprobar su consistencia.
\end{enumerate}

Esta validación demuestra que los candidatos son trayectorias factibles del
modelo completo y no únicamente patrones generados por la rutina de
búsqueda.


\subsection{Criterios de presentación}
\label{subsec:criterios_presentacion}

Los resultados se presentan de acuerdo con las siguientes reglas:

\begin{enumerate}
    \item si una cota inferior coincide con la actividad de un testigo
    factible, se reporta un mínimo global exacto;

    \item si la cota inferior es menor que el mejor testigo disponible, se
    reporta un intervalo;

    \item si la búsqueda es exhaustiva únicamente dentro de una familia,
    el mejor valor se identifica como óptimo restringido;

    \item una terminación por tiempo no se interpreta como prueba de
    infactibilidad ni de optimalidad;

    \item todo testigo empleado en las conclusiones debe ser reconstruible
    mediante los artefactos del proyecto.
\end{enumerate}

De acuerdo con estos criterios, los casos de una y dos rondas pueden
producir mínimos globales exactos. En tres rondas se obtiene una cota
inferior global y una cota superior respaldada por testigos de la familia
$2+2+c$.